%%%%%%%%%%%%%%%%%%%%%%%%%%%%%%%%%%%%%%%%%
% Medium Length Professional CV
% LaTeX Template
% Version 3.0 (December 17, 2022)
%
% This template originates from:
% https://www.LaTeXTemplates.com
%
% Author:
% Vel (vel@latextemplates.com)
%
% Original author:
% Trey Hunner (http://www.treyhunner.com/)
%
% License:
% CC BY-NC-SA 4.0 (https://creativecommons.org/licenses/by-nc-sa/4.0/)
%
%%%%%%%%%%%%%%%%%%%%%%%%%%%%%%%%%%%%%%%%%

%----------------------------------------------------------------------------------------
%	PACKAGES AND OTHER DOCUMENT CONFIGURATIONS
%----------------------------------------------------------------------------------------

\documentclass[
	%a4paper, % Uncomment for A4 paper size (default is US letter)
	11pt, % Default font size, can use 10pt, 11pt or 12pt
]{resume} % Use the resume class

% in your .tex preamble (after \documentclass)
\usepackage{tikz}

\newcommand{\linkbox}[2]{%
  \tikz[baseline=(X.base)]\node[
    inner xsep=4pt, inner ysep=2pt,
    rounded corners=2pt,
    draw=cyan!70, line width=0.6pt
  ] (X) {\href{#1}{#2}};%
}


%------------------------------------------------

\name{Yunwoo Lee} % Your name to appear at the top

% You can use the \address command up to 3 times for 3 different addresses or pieces of contact information
% Any new lines (\\) you use in the \address commands will be converted to symbols, so each address will appear as a single line.
\address{
    Phone: (206)~$\cdot$~843~$\cdot$~4857 \\ 
    Email: \href{mailto:yunu919@uw.edu}{yunu919@uw.edu} % 이메일 클릭 가능하게
} 

\address{Seattle, WA 98105}

\address{
  \linkbox{https://yunlee919.github.io/}{Website}
  \linkbox{https://www.linkedin.com/in/yunu919/}{LinkedIn}
}



%----------------------------------------------------------------------------------------

\begin{document}

%----------------------------------------------------------------------------------------
%	EDUCATION SECTION
%----------------------------------------------------------------------------------------

\begin{rSection}{Education}
	
	\textbf{University of Washington} \hfill Seattle, WA \\ 
	\textsl{Master of Science in Computational Linguistics} \hfill
    \textsl{Sep. 2025 -- Present}

    \textbf{\href{https://www.hufs.ac.kr/hufs/index.do\#section1}{Hankuk University of Foreign Studies, HUFS}} \hfill Seoul, South Korea \\
    \textsl{Bachelor of Arts in Linguistics \& Cognitive Science} \hfill
    \textsl{Mar. 2018 -- Feb. 2024} \\
	\textsl{Bachelor of Language Science in Artificial Intelligence (Double Major)}\\
    Advisor: Jeesun Nam

\end{rSection}

%----------------------------------------------------------------------------------------
%	WORK EXPERIENCE SECTION
%----------------------------------------------------------------------------------------

\begin{rSection}{Professional and Research Experience}

	\begin{rSubsection}{\href{https://www.keti.re.kr/eng/main/main.php}{Korea Electronics Technology Institute (KETI)}}{Seongnam, South Korea}{Researcher, Language Team - AIRC (Artificial Intelligence Research Center)}{Sep. 2024 -- Jun. 2025}
		\item Developed trustworthiness benchmarks for Korean LLMs (hallucination, reliability, sociocultural bias) through prompt design, rubrics, and failure-mode analysis
        \item Validated a Korean multimodal dialogue dataset by aligning utterance-level pragmatic functions with nonverbal cue labels to improve human–AI interaction robustness
	\end{rSubsection}

%------------------------------------------------

	\begin{rSubsection}{\href{https://nc-ai.com/en}{NCSOFT}}{Seongnam, South Korea}{Language AI Researcher, Language Data Team}{Mar. 2024 -- Sep. 2024}
		\item Led AI Red Team initiatives, creating robust testing protocols and diverse conversational datasets to identify and mitigate ethical and safety risks in language models
		\item Categorized language model vulnerabilities into single-turn versus multi-turn threats by identifying discourse-level risks, including anaphoric references that are undetectable within a single turn
	\end{rSubsection}

%------------------------------------------------

	\begin{rSubsection}{\href{https://www.sktelecom.com/index_en.html}{SK TELECOM}}{Seoul, South Korea}{Linguistic Annotator, AI Technology Unit}{Feb. 2023 -- Jun. 2023}
		\item Led human-centric annotation to ensure sociolinguistic coverage (honorific levels, informal slang, code-mixing) in conversational training data
		\item Created a linguistically grounded ambiguity taxonomy and applied it to normalize Korean queries and improve chatbot robustness
	\end{rSubsection}
    
%------------------------------------------------

    \begin{rSubsection}{\href{http://dicora.kr/}{DICORA Computational Linguistics Lab, HUFS}}{Seongnam, South Korea}{Undergraduate Research Intern under Prof. Jeesun Nam}{Jan. 2022 -- Sep. 2022}
		\item Contributed to two research projects (1) constructing datasets for sentiment analysis of stock-market articles, and (2) generating NLU datasets for training a big-data-driven chatbot model 
	\end{rSubsection}

\end{rSection}

%----------------------------------------------------------------------------------------
%	AWARDS AND HONORS SECTION
%----------------------------------------------------------------------------------------

\begin{rSection}{AWARDS AND HONORS}
    \textbf{\href{https://linguistics.washington.edu/honors-awards\#S06}{UW CLMS Scholarship}} (\$14,500) \hfill 2025 \\ 
    \textbf{HUFS Departmental Scholarship} \hfill 2021, 2023 \\ 
    \textbf{Outstanding Undergraduate Thesis Award} \hfill 2022 \\ 
    \textbf{Best Composition Award} \hfill 2018 \\
    \textbf{Student Leadership Scholarship} \hfill 2018 \\ 
    \textbf{National Merit Scholarship} \hfill 2018 
\end{rSection}

%----------------------------------------------------------------------------------------
%	PERSONAL PROJECT SECTION
%----------------------------------------------------------------------------------------

\begin{rSection}{PROJECT EXPERIENCE}

	\begin{rSubsection}{Linguistic Pattern Analysis for Financial Sentiment: Attribute-Verb Relationships in Stock Market Articles}{}{HUFS Linguistics Graduation Thesis Project}{Sep. 2022 -- Dec. 2022}
		\item Modeled semantic constraints between attribute nouns (e.g., interest rates, costs) and directional predicates (rise/fall) to capture context-dependent polarity shifts in financial discourse
		\item Recognized with the Outstanding Undergraduate Thesis Award
	\end{rSubsection}
\end{rSection}

%----------------------------------------------------------------------------------------
%	Language and TECHNICAL STRENGTHS SECTION
%----------------------------------------------------------------------------------------

\begin{rSection}{Languages and Technical Skills}

	\begin{tabular}{@{} >{\bfseries}l @{\hspace{6ex}} l @{}}
		Programming & Python, R, C++, SQL, Java \\
        ML/NLP & PyTorch, NLTK \\
        Tools & Git, Linux/CLI \\
        Typesetting & \LaTeX \\
		Languages & Korean (native), English (fluent) \\
	\end{tabular}

\end{rSection}

%----------------------------------------------------------------------------------------
%	EXAMPLE SECTION
%----------------------------------------------------------------------------------------

%\begin{rSection}{Section Name}

	%Section content\ldots

%\end{rSection}

%----------------------------------------------------------------------------------------

\end{document}
